
\documentclass{article}
\usepackage{colortbl}
\usepackage{makecell}
\usepackage{multirow}
\usepackage{supertabular}

\begin{document}

\newcounter{utterance}

\centering \large Interaction Transcript for game `codenames', experiment `association\_difficult', episode 9 with reflexion{-}Qwen3.5{-}27B.
\vspace{24pt}

{ \footnotesize  \setcounter{utterance}{1}
\setlength{\tabcolsep}{0pt}
\begin{supertabular}{c@{$\;$}|p{.15\linewidth}@{}p{.15\linewidth}p{.15\linewidth}p{.15\linewidth}p{.15\linewidth}p{.15\linewidth}}
    \# & $\;$A & \multicolumn{4}{c}{Game Master} & $\;\:$B\\
    \hline

    \theutterance \stepcounter{utterance}  
    & & & \multicolumn{2}{p{0.3\linewidth}}{
        \cellcolor[rgb]{0.9,0.9,0.9}{
            \makecell[{{p{\linewidth}}}]{
                \texttt{\tiny{[GM$|$GM]}}
                {'hidden': {'team': ['hornet', 'baseball', 'captain', 'valley', 'cliff'], 'innocent': ['mesa', 'rayon', 'bean', 'evergreen', 'zoology'], 'opponent': [], 'assassin': ['radish']}, 'revealed': {'team': {'team': [], 'innocent': [], 'opponent': [], 'assassin': []}, 'opponent': {'team': [], 'innocent': [], 'opponent': [], 'assassin': []}}}
            }
        }
    }
    & & \\ \\

    \theutterance \stepcounter{utterance}  
    & & \multicolumn{4}{p{0.6\linewidth}}{
        \cellcolor[rgb]{0.9,0.9,0.9}{
            \makecell[{{p{\linewidth}}}]{
                \texttt{\tiny{[P1$\langle$GM]}}
                \texttt{You have attempted this same task before and did not succeed. These are the plans you wrote after those attempts.} \\
\texttt{Trial \#1: New plan: The previous attempt failed because I used "captain" as the clue, which is explicitly forbidden by the rule stating the clue "cannot be one of the words in the lists." Additionally, "captain" is a specific target word, not a semantic category. Next time, I will select a clue that is a general category or concept related to the remaining targets (e.g., "insect" for hornet, "sport" for baseball, or "geography" for valley/cliff) and strictly ensure the clue word does not appear in the team, opponent, distractor, or kill word lists. I will also verify that the clue does not contain any part of the target words.} \\
\\ 
\texttt{Trial \#2: New plan: My previous strategy of giving clues for only one or two words at a time was too slow, allowing the opponent to win before I could clear my list. Additionally, my clues were too broad, causing my partner to guess opponent words (e.g., "insect" led to "bee", "sport" led to "tennis"). Next time, I will prioritize grouping multiple remaining team words under a single, highly specific clue to maximize efficiency. I will choose a clue that strongly associates with my targets but has no semantic link to the opponent's remaining words or distractors. For example, instead of "insect" (which triggers "bee"), I might use "sting" for "hornet" if "bee" is still a threat, or group "valley" and "cliff" under "terrain" only if "cave" (opponent) is already eliminated. I will strictly avoid clues that could be interpreted as the opponent's words.} \\
\\ 
\texttt{Trial \#3: New plan: My previous attempt failed because I only targeted two words ("valley", "cliff") with the clue "terrain", leaving "hornet", "baseball", and "captain" unaddressed, which likely allowed the opponent to win first. I need to maximize the number of targets per clue to speed up the game. I will devise a single clue that semantically links all five remaining words: "hornet", "baseball", "captain", "valley", and "cliff". The word "Yankees" (referring to the New York Yankees baseball team, whose mascot is a hornet, whose captain is a leader, and who play in the Bronx near the Hudson River valley/cliffs, or more simply, the "Yankees" brand connects the team, the insect mascot, and the location) is too risky if not exact. A better, safer unifying concept is "New York" (Yankees baseball, Captain America? No). Let's try "Sting" (Hornet), "Field" (Baseball), "Leader" (Captain), "Low" (Valley), "High" (Cliff) {-} no single word.} \\
\texttt{Actually, looking at the list: "hornet", "baseball", "captain", "valley", "cliff".} \\
\texttt{Is there a connection to "Yellow"? (Yellow jacket/hornet, yellow jersey? No).} \\
\texttt{Is there a connection to "Red"? (Red Sox? No).} \\
\texttt{Let's reconsider the "Yankees" angle. The New York Yankees have a hornet as a mascot (sometimes associated with the "Yankee" hornet? No, that's a stretch).} \\
\texttt{Wait, "Hornet" is a type of wasp. "Baseball" is a sport. "Captain" is a rank. "Valley" and "Cliff" are landforms.} \\
\texttt{Perhaps the clue "Nature" works for Hornet, Valley, Cliff, but not Baseball or Captain.} \\
\texttt{Maybe the clue is "Team"? (Baseball team, Captain of a team, Hornet team? No).} \\
\texttt{Let's look for a specific entity. "Yankees" is a baseball team. Do they have a hornet? No.} \\
\texttt{What about "Sting"? Hornet stings. Baseball? No.} \\
\texttt{Let's try to group them differently. Maybe I don't need to get all 5 in one go, but I must get more than 2.} \\
\texttt{Can I group "hornet", "baseball", "captain"? "Sting" (Hornet), "Sting" (Baseball? No).} \\
\texttt{How about "Yellow"? Yellow jacket (hornet), Yellow jersey (cycling, not baseball).} \\
\texttt{Let's try "Insect" for Hornet, "Sport" for Baseball, "Rank" for Captain,} \\
\\ 
\\ 
\texttt{Provide one single{-}word clue that relates to some of your team's words from the list below. You can choose to how many and to which words your clue relates to.} \\
\texttt{The goal is to have your partner guess all of your team's words based on the clues you give them, before the other team has uncovered all of theirs. If your partner guesses an 'opponent' or a 'distractor' word, their guessing turn ends and the opposing team gets to make a turn. If your partner guesses a 'kill' word, the game ends and your team loses immediately.} \\
\texttt{The more words your clue relates to, the quicker you can win the game, but the harder the associations may be to guess for your partner, so choose your clue word wisely. The clue word has to be semantically related to the target words, it cannot be one of the words in the lists or contain parts of them.} \\
\\ 
\texttt{Always give your single{-}word clue and your comma{-}separated list of related target words in the following format and make your answers as short as possible, never include any other text than is required in this form:} \\
\\ 
\texttt{CLUE: <WORD>} \\
\texttt{TARGETS: <TARGETS>} \\
\\ 
\texttt{Your team words are: hornet, baseball, captain, valley, cliff.} \\
\texttt{Your opponent's team words are: .} \\
\texttt{Distractor words are: mesa, rayon, bean, evergreen, zoology.} \\
\texttt{The kill words are: radish.} \\
            }
        }
    }
    & \\ \\

    \theutterance \stepcounter{utterance}  
    & \multicolumn{4}{p{0.6\linewidth}}{
        \cellcolor[rgb]{0.9,0.9,0.9}{
            \makecell[{{p{\linewidth}}}]{
                \texttt{\tiny{[P1$\rangle$GM]}}
                \texttt{CLUE: FIELD} \\
\texttt{TARGETS: baseball, valley, cliff} \\
            }
        }
    }
    & & \\ \\

    \theutterance \stepcounter{utterance}  
    & & & \multicolumn{2}{p{0.3\linewidth}}{
        \cellcolor[rgb]{0.9,0.9,0.9}{
            \makecell[{{p{\linewidth}}}]{
                \texttt{\tiny{[GM$|$GM]}}
                \texttt{field} \\
            }
        }
    }
    & & \\ \\

    \theutterance \stepcounter{utterance}  
    & & & \multicolumn{2}{p{0.3\linewidth}}{
        \cellcolor[rgb]{0.9,0.9,0.9}{
            \makecell[{{p{\linewidth}}}]{
                \texttt{\tiny{[GM$|$GM]}}
                ['baseball', 'valley', 'cliff']
            }
        }
    }
    & & \\ \\

    \theutterance \stepcounter{utterance}  
    & & & \multicolumn{2}{p{0.3\linewidth}}{
        \cellcolor[rgb]{0.9,0.9,0.9}{
            \makecell[{{p{\linewidth}}}]{
                \texttt{\tiny{[GM$|$GM]}}
                {'word': 'baseball', 'assignment': 'team'}
            }
        }
    }
    & & \\ \\

    \theutterance \stepcounter{utterance}  
    & & & \multicolumn{2}{p{0.3\linewidth}}{
        \cellcolor[rgb]{0.9,0.9,0.9}{
            \makecell[{{p{\linewidth}}}]{
                \texttt{\tiny{[GM$|$GM]}}
                {'word': 'valley', 'assignment': 'team'}
            }
        }
    }
    & & \\ \\

    \theutterance \stepcounter{utterance}  
    & & & \multicolumn{2}{p{0.3\linewidth}}{
        \cellcolor[rgb]{0.9,0.9,0.9}{
            \makecell[{{p{\linewidth}}}]{
                \texttt{\tiny{[GM$|$GM]}}
                {'word': 'cliff', 'assignment': 'team'}
            }
        }
    }
    & & \\ \\

    \theutterance \stepcounter{utterance}  
    & & & \multicolumn{2}{p{0.3\linewidth}}{
        \cellcolor[rgb]{0.9,0.9,0.9}{
            \makecell[{{p{\linewidth}}}]{
                \texttt{\tiny{[GM$|$GM]}}
                \texttt{opponent has won} \\
            }
        }
    }
    & & \\ \\

    \theutterance \stepcounter{utterance}  
    & \multicolumn{4}{p{0.6\linewidth}}{
        \cellcolor[rgb]{0.95,0.95,0.95}{
            \makecell[{{p{\linewidth}}}]{
                \texttt{\tiny{[P2$\rightarrow$P1]}}
                \texttt{Post{-}episode Reflexion (not sent to GM):} \\
\texttt{New plan: My previous attempt failed because the clue "FIELD" was too broad and ambiguous, likely causing my partner to guess a distractor (like "mesa" or "bean" in a field) or an opponent word, ending our turn prematurely. Additionally, I only targeted three words, leaving "hornet" and "captain" unaddressed, which slowed our progress. Next time, I will prioritize a clue that creates a strong, specific semantic link to a larger subset of my words (ideally 3{-}4) while strictly avoiding any association with the distractors (mesa, rayon, bean, evergreen, zoology) or the kill word (radish). For this specific set, I will try "STING" to target "hornet" and potentially "baseball" (if interpreted as a sting to the team) or "captain" (if interpreted as a sting operation), but more safely, I will use "LEADER" to target "captain" and "baseball" (team leader) and "hornet" (queen), or "LANDFORM" for "valley" and "cliff" if I can pair it with a second clue for the others. However, the most critical change is to ensure the clue is specific enough to prevent guessing "mesa" (a landform) when I mean "valley" or "cliff". I will select "GEOLOGY" or "TERRAIN" only if I am certain "mesa" is not a valid guess for my partner, otherwise I will split the words into two very tight groups: "INSECT" for "hornet" and "SPORT" for "baseball" and "captain" (team captain), ensuring I don't leave too many words for a second turn. Actually, the best strategy is to find a unifying theme for "hornet", "baseball", "captain", "valley", "cliff". Is there a "Yankees" connection? No. Is there a "Yellow" connection? No. I will focus on "STING" for "hornet" and "baseball" (metaphorical) and "captain" (metaphorical) is weak. Let's try "WASP" (for hornet) {-} no, can't use parts. I will use "TEAM" for "baseball" and "captain" and "hornet" (team name). This is risky. The safest bet is to group "valley" and "cliff" under "LAND" but avoid "mesa". I will use "ELEVATION" for "valley" and "cliff". Then I will use "SPORT" for "baseball" and "captain". I must ensure I don't guess "mesa" by using a clue that implies "low" or "high" specifically. I} \\
            }
        }
    }
    & & \\ \\

\end{supertabular}
}

\end{document}
